\section{Discussion and limitations}
\label{sec:discussion}

% Skeleton with draft paragraphs, written before the results. The limitations are written first, from the
% preregistration ("What is not concluded") and Section 3.5; the discussion of the results is added after them.

\subsection{What the results do not show}
\label{sec:discussion:limits}

\paragraph{No proof of safety.}
The results describe the attackers tried here, under the assumptions of Section~\ref{sec:threat-model:bias}. They give
no proof of safety and no upper bound on the risk of matching, and a form that defeats these attackers may not defeat a
stronger one. The strongest attacker is the maximum over the attackers tried, not an attack that combines all the
length and timing information a form leaves; only A3 on R3-u combines the two. The alignment for suppression ignores
the binomial variation and is not a likelihood. The timing term of A3 on R3-u uses, in effect, only the part of a log
before the first error of the replay.

\paragraph{Paths that the forms leave open.}
In the forms that release press times, the duration of a drag, and that of a held key when the next press follows its
release, remains as a lower bound of the next press interval. No attacker targets it, but it can enter the interval
features of A2-b and the segment durations of A3 on R7. The rule that treats the next press after Control as a shortcut
fails when several keys are pressed while Control is held, because whether Control is held cannot be read when its
action time is zero.

\paragraph{The closed world and the candidates.}
The true text is always among the candidates, the attackers always choose one and never abstain, and the open-world
diagnostic does not give a rate of false matches. The decoys are drawn uniformly from the writers of the same prompt,
who may be more or less alike than the members of a class. On windows the attacker knows both ends of the window. The
distribution of accuracy per writer is reported only up to its 90th percentile.

\paragraph{What the data allow.}
We claim no difference by writer attributes, since the demographic file was not acquired, and we do not describe matching
as identification of writers. The results on KLiCKe are not generalised to other tasks (a limit of 30 minutes,
argumentative essays, crowdsourced writers) or to logs that did not pass the selection of the KLiCKe release, such as
logs with much pasting~\cite{tian2025klicke}. The results on ENJA15 are not generalised to Japanese free writing or to
logs of input methods in general, and differences from KLiCKe are not attributed to language or input method, since
the task, the input and the language differ at once. Our formulas for the process measures could not be verified
against Inputlog (\S\ref{sec:protocol:utility}). KeyTAR was run with the defaults of its artifact, including a limit on
generated tokens that is expected to cut reconstructions short, and is a reference, not a bound.

\paragraph{The exploratory analyses.}
The uniqueness of word-length sequences is measured among the $N$ candidates tried and is not a risk of
re-identification in a population. The form with differentially private noise gives its guarantee for one segment
count or one interval, not for a writer's log, and its noise comes from a seeded generator for evaluation, not from an
implementation suitable for release.

\subsection{Processing and agreements}
\label{sec:discussion:agreements}

Processing is not the only protection that a releasing party can give. A release can also be bound by terms of
use. The terms of DataShop, a repository of data from educational software, state ``You will not make any attempts to
de-anonymize the data'' and forbid creating a dataset with information that identifies
students~\cite{datashop2012terms}. The rules of the competition allow the data to be used for non-commercial purposes,
including academic research, and forbid passing them to anyone who has not agreed to the rules; in the rules and on the
data page we found no clause on recovering the texts or identifying the writers~\cite{kaggle2023lwp}. KLiCKe may be
downloaded and used freely, with the repository cited as the source~\cite{inputlog2026community}. Terms do not change
what a log reveals to a party that ignores them, and they do not bind a party that holds the texts and obtains the
logs outside the terms; but they define who may use a release and for what, and they can be combined with any of the
forms evaluated here.

\subsection{What a releasing party can take away}
\label{sec:discussion:takeaways}

\textbf{[After the results.]} The results are material for three choices: which columns to keep, how coarsely to give
times and lengths, and whether to release full logs or windows. Each statement takes the form ``against this attacker,
at this $N$, this form had this multiple of $1/N$ and kept this much of the process measures''. No rule recommends a
form, and the paper does not say that a form works or withstands matching.
